\section{The Unease of Renormalization}

\label{sec:renorm_clue}

Renormalization is one of the most successful procedures in modern physics.
Quantum field theory produces divergent quantities during calculation;
renormalization absorbs these divergences into observable definitions
and yields predictions that agree with experiment to extraordinary precision.

Yet this success conceals a fundamental unease.

The procedure works with remarkable accuracy,
yet its conceptual meaning remains unresolved.
A theory that produces infinities mid-calculation
is nevertheless able to describe physical reality correctly
after appropriate redefinition.

This produces a peculiar situation.

A theory breaks down once,
yet after modification,
it succeeds.

Why does a theory that diverges still succeed?

This question is rarely pursued at a foundational level.
Renormalization has been formalized and refined,
and its mathematical structure is well understood.
However, its physical meaning remains unclear.

What, precisely, is renormalization doing?

Why do divergences appear?
And why does removing them yield correct predictions?

No unified interpretation of these questions has been established.

In this paper, we take this unease seriously.

Rather than treating renormalization as a problem to be resolved,
we treat it as a structural clue.

The appearance of divergence may not indicate a failure of theory,
but contact with the limit of its descriptive window.

From this perspective,
renormalization is not merely a technical operation.
It may be a procedure that converts boundary contact
into finite relations that can be handled within the theory.

This reinterpretation will be developed throughout this volume.

Divergence is not removed;
it is transformed.

And what appears as infinity
may be the signal of a descriptive boundary.
